% ED26-SCI-Mahajan2014-002; selected Mahajan 2014 Sections 6.1--6.4, 8.1--8.4.
\chapter{Approximate Quantitative Inference}
\label{sec:p15-approximate-quantitative-inference}

An approximation is useful when it preserves the conclusion that matters.
Estimating a quantity within ten percent can settle a factor-of-two comparison
and still fail to settle its sign after subtraction from a nearby threshold.
This chapter develops Mahajan's methods of logarithmic rounding, characteristic
values, shape lumping and easy cases
\cite[Secs.~6.1--6.4, 8.1--8.4]{Mahajan2014Insight}.
The examples retain their model assumptions. Additional remainder and decision
bounds are proved here; they are not attributed as the source's theorems.
Dimensionless formulation and scaling are treated in
\cref{sec:p15-dimensional-and-scaling-inference}.

\section{State the relation before using the number}

\begin{hbdefinition}{def:p15-approximation-relations}{Four different approximation claims}
An \emph{identity} \(f(x)=g(x)\) asserts equality throughout a declared domain.
A \emph{bound} specifies an inequality, such as
\(|f(x)-g(x)|\leq E(x)\), throughout that domain.
For a parameter \(\varepsilon\to0^+\), with all other variables fixed,
\(f\sim g\) means \(f/g\to1\), where \(g\) is eventually nonzero.
The assertion \(f=g+O(h)\) means that some constants \(C,\varepsilon_0>0\)
satisfy \(|f-g|\leq C|h|\) for \(0<\varepsilon<\varepsilon_0\).
It does not supply numerical values for those constants.

A \emph{heuristic estimate} replaces a process or parameter by a convenient
model value without supplying such a guarantee. Mahajan grades informal
equalities by what they may omit
\cite[Sec.~1.1, p.~5]{Mahajan2014Insight}. Besides \(\equiv\) (equality by
definition) and \(=\), he writes \(f\approx g\) when the two sides may differ
by a purely numerical factor near \(1\), \(f\sim g\), read ``is comparable
to'', when they may differ by a purely numerical factor of any size, and
\(f\propto g\) when the omitted factor may carry dimensions. For a circle of
radius \(r\) and area \(A\) his examples are \(A=\pi r^2\), \(A\approx4r^2\)
and \(A\sim r^2\), with omitted factors \(\pi/4\) and \(\pi\). For a cylinder
of height \(h\), \(V\sim hr^2\) gives \(V\propto r^2\) and \(V\propto h\);
the factor omitted in \(V\propto h\) has dimensions of length squared.
None of the three relations is the ratio limit above. The drag force is
written \(F_{\mathrm d}\sim\rho v^2A_{\mathrm{cs}}\)
\cite[Sec.~8.3, p.~297, eq.~(8.33)]{Mahajan2014Insight}. The drag
coefficient there is
\[
 c_{\mathrm d}=\frac{F_{\mathrm d}}{\frac12\rho v^2A_{\mathrm{cs}}},
\]
so the omitted factor is \(c_{\mathrm d}/2\). It equals \(0.2\), \(0.25\),
\(0.5\) and \(1\) for the four shapes tabulated on that page. No limit is
taken, and the factor need not approach \(1\).

This chapter reserves \(\sim\) for ratio limits and writes
\(f\approx_{\mathrm{est}}g\) for the source's \(f\sim g\): a heuristic claim
that \(f/g\) is a dimensionless number of unstated size. Each
\(\approx_{\mathrm{est}}\) below restates a \(\sim\) of the source. Where the
source's stronger \(\approx\) is used, the text states that the omitted factor
is claimed to be near \(1\), as for the terminal speed
\(v_{\mathrm{term}}\approx1\,\mathrm{m\,s^{-1}}\) of the falling cone
\cite[Sec.~6.4, p.~218]{Mahajan2014Insight}. The relation \(\propto\) keeps
the source's meaning, with the held-fixed quantities stated as in
\cref{def:p15-dimensional-scaling-law}. Neither \(\approx\) nor \(\sim\) says
how near \(1\), or how large, the omitted factor is; the claim becomes
checkable only when a factor bound is stated.
Positive quantities satisfy a stated \emph{factor bound} \(K\geq1\) when
\(K^{-1}\leq f/g\leq K\). Numerical rounding error, model error and statistical
uncertainty are separate contributions; one does not bound another.
\index{approximation!meaning}\index{asymptotic equivalence}\index{lumping}
\end{hbdefinition}

For example, \cref{thm:p15-stirling-with-remainder} gives an explicit finite
remainder as well as Stirling's asymptotic relation. The uniform assumptions
in \cref{thm:p15-laplace-relative-error} must be checked before applying
that particular theorem. Neither result permits an arbitrary finite input to be
declared accurate merely because a limiting approximation exists.

\begin{hbtheorem}{thm:p15-approximation-certificates}{proved=annex:p15.approximation.certificates}{A bound that preserves a decision}
Suppose \(L\leq f\leq U\). Then \(L>c\) certifies \(f>c\), and \(U<c\)
certifies \(f<c\). If \(c\in[L,U]\), this interval alone certifies neither
strict comparison. If \(|f-\widehat f|\leq E_f\) and
\(|g-\widehat g|\leq E_g\), then
\[
 |(f-g)-(\widehat f-\widehat g)|\leq E_f+E_g.
\]
Thus \(\widehat f-\widehat g>E_f+E_g\) certifies \(f>g\).
If \(|\widehat f|>E_f\), the relative-error bound is
\(E_f/(|\widehat f|-E_f)\). No relative error is defined at a zero target.
\end{hbtheorem}

\section{Rounding on a logarithmic scale}

\begin{hbdefinition}{def:p15-approximation-log-round}{Logarithmic rounding}
For \(x>0\), base \(b>1\) and exponent spacing \(h>0\), choose an integer
\(j\) nearest to \(\log_b(x)/h\), resolving an exact tie by either fixed rule,
and set \(Q_h(x)=b^{hj}\). This rounds ratios rather than differences.
The boundary between consecutive representatives is their geometric mean.
The base and exponent act on a numerical value in stated units; a logarithm
of a physical quantity requires division by a reference quantity of the same
dimension.
\end{hbdefinition}

\begin{hbtheorem}{thm:p15-approximation-log-bound}{proved=annex:p15.approximation.log-bound}{Exact factor budget for logarithmic rounding}
The rounding satisfies
\[
 b^{-h/2}\leq Q_h(x)/x\leq b^{h/2}.
\]
For positive inputs \(x_i\) and fixed real powers \(a_i\),
\[
 b^{-\frac h2\sum_i|a_i|}
 \leq\frac{\prod_i Q_h(x_i)^{a_i}}{\prod_i x_i^{a_i}}
 \leq b^{\frac h2\sum_i|a_i|}.
\]
No independence of the rounding errors is assumed. In particular, rounding
each factor independently does not guarantee cancellation of their errors.
\end{hbtheorem}

\begin{hbexample}{def:p15-approximation-calendar}{Minutes, seconds and a mnemonic}
Nearest integral powers of ten replace \(24\) by \(10\) and \(60\) by
\(100\). Their product estimates the \(1440\) minutes in a 24-hour day by
\(1000\), an error of \(440/1440=11/36\) of the exact target.
The rounding boundary between \(10\) and \(100\) is \(\sqrt{1000}\), not
their arithmetic midpoint \(55\).

Using half powers instead gives
\(365\mapsto10^{2.5}\), \(24\mapsto10^{1.5}\), and
\(3600\mapsto10^{3.5}\). Hence a stipulated 365-day year is estimated by
\(10^{7.5}\) seconds, compared with its exact \(31\,536\,000\) seconds.
The convenient mnemonic \(\pi\times10^7\) is a further numerical
approximation, about \(0.381\) percent low. Its \(\pi\) comes from the
chosen mnemonic and units, not from a proof about circular or elliptical
orbits. With the source's circular-orbit model and radius
\(1.5\times10^{11}\,\mathrm m\), the speed estimate
\(2\pi r/(\pi10^7\,\mathrm s)=3\times10^4\,\mathrm{m/s}\) follows.
\end{hbexample}

Mahajan motivates this method with historical multiplication-estimation
responses and approximate visual discrimination. Those examples motivate
a reasoning exercise; they do not establish a universal perceptual threshold
or prove that less accurate arithmetic improves every decision.

\section{What a characteristic value retains}

\begin{hbtheorem}{thm:p15-approximation-lumping}{proved=annex:p15.approximation.lumping}{Linear lumping and its nonlinear remainder}
Let \(w_i\geq0\), \(\sum_iw_i=1\), and \(\mu=\sum_iw_ix_i\).
For an affine function \(a+bx\), replacing the inputs by \(\mu\) preserves
their weighted output exactly. If \(f\) is twice continuously differentiable
on an interval containing the \(x_i\) and \(\mu\), with \(|f''|\leq M\),
then
\[
 \left|\sum_iw_if(x_i)-f(\mu)\right|
 \leq\frac M2\sum_iw_i(x_i-\mu)^2.
\]
For the square function the discrepancy is exactly
\(\sum_iw_ix_i^2-\mu^2=\sum_iw_i(x_i-\mu)^2\).
\end{hbtheorem}

\begin{hbexample}{def:p15-approximation-populations}{Representative regions and rectangular age profiles}
Mahajan's historical example replaces 50 unequal state populations by
\(6\) million each, giving \(300\) million, and compares it with the
reported 2012 population of \(314\) million. The exact identity is
\(\sum_i n_i=50\overline n\); selecting one familiar state as the average
is the heuristic step. If \(|n_i-n_*|\leq d_i\), the total error is at most
\(\sum_i d_i\). A uniformly sampled state gives an unbiased total estimate,
but that does not make one realized estimate accurate.

The source also replaces an age-density curve by a rectangle of area
\(N=300\) million and width \(L=75\) years. Its height is
\(N/L=4\) million per year. Four age years and a stipulated participation
fraction \(q=1/2\) give \(q(4N/L)=8\) million students. This is a
steady-population and participation model, not a census identity. Birth-rate
changes, mortality, migration, study duration and older students can all
change its result.

With the source's separate 1950 UK inputs, \(N=50\) million and \(L=65\)
years, the same participation assumption predicts approximately
\(385\,000\) graduates per year. A reported \(17\,000\) would instead
require \(q\approx0.0221\) within this particular rectangular model.
Participation percentages with another denominator cannot be substituted
into that inversion without changing the model.
\end{hbexample}

\begin{hbtheorem}{thm:p15-approximation-energy-average}{proved=annex:p15.approximation.energy-average}{The energy lost by replacing velocities by one velocity}
For masses \(m_i>0\), velocities \(v_i\), total mass \(M_0=\sum_i m_i\)
and common final velocity \(\overline v=\sum_i m_iv_i/M_0\),
\[
 \frac12\sum_i m_iv_i^2-\frac12M_0\overline v^2
 =\frac12\sum_i m_i(v_i-\overline v)^2\geq0.
\]
Equality holds exactly when every velocity is the same. Thus equalizing
velocities conserves momentum but generally does not conserve kinetic energy.
\end{hbtheorem}

The source's two unit-mass slabs at velocities \(6\) and \(4\) end at
\(5\), with kinetic energy falling from \(26\) to \(25\). A physical model
that converts this difference into heat needs a dissipative mechanism; the
algebra identifies the available energy difference.

\begin{hbexample}{def:p15-approximation-jumps}{A conditional size-independent jump height}
Let \(m\) be body mass, \(\alpha m\) active muscle mass, \(e\) supplied
energy per muscle mass, and \(\eta\) the fraction converted into upward
gravitational potential energy. For a standing jump with no other energy
input and no accounted drag loss, the model balance is
\[
 mgh=\eta\alpha m e,\qquad h=\eta\alpha e/g.
\]
Mass cancels only while \(\eta,\alpha,e,g\) are held fixed. The source's
four reported examples are a \(0.5\)-mg flea reaching \(20\) cm, a
\(40\)-mg click beetle reaching \(30\) cm, a \(3\)-g locust reaching
\(59\) cm, and a \(70\)-kg human reaching \(60\) cm. This small historical set illustrates
a scaling hypothesis; it is not a proof that all animals share these parameters.
For geometrically similar jumpers, fixed target height and approximately
quadratic drag, drag work divided by \(mgh\) scales as area/mass,
approximately the inverse linear size. It can therefore become important
at the small-size end where the initial model omits it.
\end{hbexample}

\begin{hbtheorem}{thm:p15-approximation-oscillators}{proved=annex:p15.approximation.oscillators}{Characteristic times with a specified equation}
The model \(m\ddot x+kx=0\), with \(m,k>0\), has nonconstant periodic
solutions of period \(2\pi\sqrt{m/k}\), independent of amplitude.
More generally, for the restoring force
\(m\ddot x=-k|x|^{n-1}x\), with \(n>0\) and force defined as zero at
\(x=0\), a motion of amplitude \(A>0\)
has period
\[
 T(A)=4\sqrt{\frac{m(n+1)}{2k}}\,
 A^{(1-n)/2}\int_0^1\frac{du}{\sqrt{1-u^{n+1}}}.
\]
The integral is finite. The characteristic-time estimate retains the power
\(A^{(1-n)/2}\); it does not determine the prefactor by dimensional
cancellation alone.
\end{hbtheorem}

For a small-angle ideal pendulum, the linearized equation is
\(\ddot\theta+(g/l)\theta=0\), giving \(2\pi\sqrt{l/g}\) within that
model. Large angles require a different approximation check. Likewise,
replacing a derivative by characteristic change divided by characteristic
time is a scale estimate. In uniform circular motion it must use the
velocity vector, not the constant speed: differentiation of
\(r(\cos\omega t,\sin\omega t)\) gives acceleration magnitude
\(r\omega^2=v^2/r\).

\section{Lumping a shape or a changing process}

\begin{hbtheorem}{thm:p15-approximation-rectangle}{proved=annex:p15.approximation.rectangle}{An error budget for replacing a graph by a rectangle}
On an interval \(I\) of finite length \(L\), if \(f\) is integrable and
\(|f(t)-c|\leq d\), then
\[
 \left|\int_I f(t)\,dt-cL\right|\leq dL.
\]
For an integrable function on a larger domain, an omitted region contributes
at most its integral of \(|f|\). If \(f\) is Lipschitz with constant \(K\)
on \([t_0-L/2,t_0+L/2]\), using height \(f(t_0)\) gives error at most
\(KL^2/4\). A peak height and a guessed width without these bounds give
a heuristic rectangle, not a certified integral.
\end{hbtheorem}

\begin{hbexample}{def:p15-approximation-shape-models}{Cells, molecular paths and a localized field}
Approximating an atom by a cube of side \(a\) gives the cell model
\(\rho=A m_p/a^3\). Mahajan uses \(a\approx_{\mathrm{est}}3\) angstroms and
\(m_pN_A\approx1\,\mathrm{g/mol}\), obtaining roughly
\(\rho\approx_{\mathrm{est}}A/18\,\mathrm{g/cm^3}\).
If the true spacing is a factor \(s\) times the assumed spacing, the density
changes by \(s^{-3}\); a density discrepancy of \(2.5\) corresponds to a
spacing smaller by \(2.5^{1/3}\). Molecular composition, packing and voids
must be included before treating this as a material-specific prediction.

For a moving sphere of radius \(r\) among fixed point centers of number
density \(n\), the swept-tube estimate uses area \(\sigma=\pi r^2\) and
\(n\sigma\lambda\approx_{\mathrm{est}}1\). Equal-radius targets replace
\(\sigma\) by \(\pi(2r)^2\). Under the additional model that centers form
a homogeneous Poisson process, the no-collision probability over distance
\(s\) is \(e^{-n\sigma s}\), and its integral gives the exact model mean
\(\lambda=1/(n\sigma)\). Moving correlated molecules require another model.
The source's \(22\)-liter molar volume and \(3\)-angstrom diameter give an
order-of-magnitude path of \(10^{-7}\,\mathrm m\).

A localized weak-deflection estimate takes acceleration scale \(GM/b^2\)
and traversal time \(b/c\), giving a transverse-speed scale \(GM/(bc)\)
and angle scale \(GM/(bc^2)\). The omitted integral and the gravitational
model determine the coefficient. Extrapolating this weak-angle estimate to
order-one deflection does not prove capture or the existence of an event
horizon. It signals that the weak-deflection approximation needs replacement.
\end{hbexample}

For a generic impulsive energy release into a uniform medium, the source's
characteristic-value argument uses energy density scale \(E/R^3\), speed
scale \((E/(\rho R^3))^{1/2}\), and \(R/t\) of comparable size. Its
consequence is the scaling \(E\approx_{\mathrm{est}}\rho R^5/t^2\).
It concerns a stipulated self-similar strong-shock regime; a numerical
coefficient, energy losses, background pressure and departures from that
regime are not determined by the scaling argument.

\begin{hbtheorem}{thm:p15-approximation-fall}{proved=annex:p15.approximation.fall}{A check on the terminal-speed approximation}
In the model \(\dot v=g(1-v^2/V^2)\), \(v(0)=0\), with \(g,V>0\),
\[
 v(t)=V\tanh(gt/V),\qquad
 s(t)=\frac{V^2}{g}\log\cosh(gt/V).
\]
For \(t\geq0\),
\[
 0\leq Vt-s(t)\leq\frac{V^2}{g}\log2,\qquad
 0\leq1-v(t)/V\leq2e^{-2gt/V}.
\]
At a prescribed distance \(s\),
\(v/V=\sqrt{1-e^{-2gs/V^2}}\).
\end{hbtheorem}

The source's lumped cone model uses \(V=1\,\mathrm{m/s}\) and
\(g=10\,\mathrm{m/s^2}\). Its free-fall-to-terminal corner occurs at
\(0.1\) s and \(5\) cm. The quadratic-drag model instead gives speed
fraction \(\tanh1\approx0.762\) at \(0.1\) s, and
\(\sqrt{1-e^{-1}}\approx0.795\) at \(5\) cm. Thus ``a substantial fraction''
is meaningful here; exact arrival at terminal speed takes no finite time.
Over \(2\) s the constant-terminal distance error is at most about
\(6.94\) cm in this model.

\begin{hbexample}{def:p15-approximation-contact}{A localized elastic contact model}
For a sphere of radius \(R\), small indentation \(\delta\) and contact
radius \(r\), the geometry gives \(r^2=2R\delta-\delta^2\).
When \(\delta/R\ll1\), this gives \(\delta\approx r^2/(2R)\).
The source's lumped strain region has depth comparable to \(r\), not \(R\);
stress scale \(Y\delta/r\) acting over area scale \(r^2\) gives restoring
force scale \(Yr^3/R\). Balancing a weight scale \(\rho gR^3\) yields
\[
 r/R\approx_{\mathrm{est}}(\rho gR/Y)^{1/3}.
\]
The result is self-consistent only when this ratio is small, the response
is elastic and the localized strain model applies. The source's rubber-ball
values \(R=10^{-2}\,\mathrm m\), \(\rho=10^3\,\mathrm{kg/m^3}\),
\(Y=3\times10^7\,\mathrm{Pa}\) give a contact-radius scale near
\(10^{-4}\,\mathrm m\). This estimate has no certified material error bar.
Using a specified contact law supplies the omitted coefficient and a
testable validity condition.
\end{hbexample}

\section{Easy cases reject formulas; they do not prove them}

\begin{hbtheorem}{thm:p15-approximation-easy-cases}{proved=annex:p15.approximation.easy-cases}{Finite checks and matching limits are insufficient}
A finite collection of values on an interval containing other points does
not determine a function without a restricting model class. A candidate \(f\) can be changed to
\(f+c\prod_i(x-x_i)^2\), preserving every checked value \(x_i\).
Likewise, for any \(c>0\),
\[
 g_c(x)=\frac1{1+x}\left(1+\frac{cx}{(1+x)^2}\right),\quad x\geq0,
\]
satisfies \(g_c(0)=1\) and \(g_c(x)\sim1/x\) as \(x\to\infty\), yet
\(g_c(1)=(1+c/4)/2\). Matching both limiting regimes supplies no uniform
accuracy guarantee in the transition.
\end{hbtheorem}

\begin{hbtheorem}{thm:p15-approximation-counting-cases}{proved=annex:p15.approximation.counting-cases}{Two exact counting checks}
Inclusive integer indices \(m_1\leq m_2\) number \(m_2-m_1+1\).
For integers \(d\geq1,n\geq0\), with \(n\) independent uniform
assignments to \(d\) possible birthdays,
\[
 P(\text{some match})=1-\prod_{j=0}^{n-1}(1-j/d),\quad 0\leq n\leq d,
\]
with empty product one; for \(n>d\), a match is certain.
\end{hbtheorem}

The source's hypothetical \$500 divided among August through December
therefore gives \$100 per month. December alone exposes the erroneous
denominator \(m_2-m_1\). Similarly, one person's match probability is zero,
which rejects the birthday product with one extra factor. The counting
derivation, rather than that one check alone, proves the surviving formula.

\begin{hbtheorem}{thm:p15-approximation-mechanical-cases}{proved=annex:p15.approximation.mechanical-cases}{Models that determine the missing functions}
Under the models specified below, with \(g>0\) and launch speed \(v\geq0\):
\begin{enumerate}
\item An ellipse with semiaxes \(a,b>0\) has area \(\pi ab\).
\item An ideal Atwood machine with masses \(m_1,m_2>0\), massless
inextensible string and massless frictionless pulley has downward
acceleration of mass 1 and string tension
\[
 a=g\frac{m_1-m_2}{m_1+m_2},\qquad
 T=\frac{2gm_1m_2}{m_1+m_2}.
\]
\item A projectile launched and received at the same height in constant
gravity, without drag, has horizontal range
\(R=2v^2\sin\theta\cos\theta/g\), for \(0\leq\theta\leq\pi/2\).
\item A rigid body rolling without slipping with fixed rolling radius
\(r>0\), mass \(m>0\), moment of inertia \(I\geq0\) and adequate static
friction has acceleration
\(a=g\sin\theta/(1+I/(mr^2))\) down a fixed incline.
\end{enumerate}
Each formula needs the corresponding geometry or dynamical equations;
symmetry and limiting checks alone do not determine it.
\end{hbtheorem}

The Atwood parameter \((m_1-m_2)/(m_1+m_2)\) changes sign under exchange of
the masses, just as the signed acceleration does. In the rolling example,
\(I/(mr^2)\) can exceed one when a narrow axle carries distant masses:
\(r\) is the rolling radius, not the largest radius of the object. A disc
and ring with the usual mass distributions have parameters \(1/2\) and
\(1\); their accelerations are respectively \(2g\sin\theta/3\) and
\(g\sin\theta/2\).

\section{Dominance, competition and the transition region}

\begin{hbtheorem}{thm:p15-approximation-dominance}{proved=annex:p15.approximation.dominance}{A quantitative dominant-term criterion}
For \(A,B>0\), let \(M=\max(A,B)\), \(m=\min(A,B)\), and \(r=m/M\).
Then \(M\leq A+B\leq2M\), and using \(M\) has relative error
\(r/(1+r)\) against the true sum. Separately, for real \(A,B\) with
\(A\ne0\), \(|B|\leq r|A|\) and \(0\leq r<1\),
then omitting \(B\) from \(A+B\) has relative error at most \(r/(1-r)\).
The denominator records the danger of cancellation.
\end{hbtheorem}

\begin{hbexample}{def:p15-approximation-drag}{Choose a drag regime and check it afterward}
Let \(\nu\) be kinematic viscosity, \(\rho\) fluid density, \(r\) sphere
radius and \(v\) relative speed. The diameter-based Reynolds number is
\(\mathrm{Re}=2rv/\nu\). Under the stipulated creeping-flow sphere law
\(F_d=6\pi\rho\nu rv\), its drag coefficient in
\(F_d=\tfrac12\rho\pi r^2v^2C_d\) is \(C_d=24/\mathrm{Re}\).
The source contrasts this with an approximately constant high-Reynolds
coefficient. The latter is a regime-dependent empirical model; taking
viscosity literally to zero is not a proof of its value or of a uniform
inviscid limit. The source's illustrative high-Reynolds coefficients are
about \(0.4\) for a car, \(0.5\) for a sphere, \(1\) for a cylinder and
\(2\) for a transverse flat plate. They are not shape-independent constants.
Its sphere example also identifies a marked coefficient change near
Reynolds number \(3\times10^5\), so an approximately constant region
must not be extrapolated through every later transition.

For a dense falling sphere, neglecting buoyancy gives the creeping-flow
terminal-speed candidate \(v=2\rho_sgr^2/(9\rho\nu)\).
Substitute this result into \(2rv/\nu\). If that value is not small, the
calculation has contradicted its own premise. Using illustrative air values
\(\rho=1.2\), \(\nu=1.5\times10^{-5}\) in SI units and
\(\rho_s=1000\), a \(10\)-micrometer radius gives about
\(0.012\,\mathrm{m/s}\) and \(\mathrm{Re}\approx0.016\).
A \(2.5\)-millimeter radius gives a creeping-flow candidate hundreds of
meters per second and a Reynolds number above \(10^5\), decisively failing
the premise. The contradiction selects further modeling work, not that
unphysical speed as a prediction.
\end{hbexample}

\begin{hbtheorem}{thm:p15-approximation-competing-energies}{proved=annex:p15.approximation.competing-energies}{Balancing scales and locating an actual minimum}
For \(A,B>0\), the stipulated energy \(E(r)=A/r^2-B/r\), \(r>0\), has
unique minimum at \(r_0=2A/B\), with \(E_0=-B^2/(4A)\).
Equating the magnitudes of its two terms instead gives \(A/B\); it
predicts the scale, not the exact minimizer.
At energy \(E_0+\Delta\), \(0<\Delta<|E_0|\), the two turning radii have
midpoint
\[
 \frac{r_-+r_+}{2}=\frac{r_0}{1-\Delta/|E_0|}>r_0.
\]
This midpoint is not by definition a time average or a thermal expectation.
\end{hbtheorem}

The source uses competition between confinement and attraction to estimate
an atomic size, and a skew potential to explain a tendency toward thermal
expansion. The theorem gives a precise version for one stipulated potential.
An actual thermal expansion coefficient also depends on the potential,
statistical distribution, constraints and phase; a sketch of asymmetric
turning points does not prove a universal positive coefficient.

The atmospheric example compares \(mgH\) with \(k_BT\). In the ideal,
isothermal, constant-gravity model, hydrostatic balance
\(dP/dz=-\rho g\), with \(P=nk_BT\) and \(\rho=mn\), gives
\(n(z)=n(0)e^{-z/H}\), \(H=k_BT/(mg)\). This is a scale height, not a
sharp top of the atmosphere. The stellar example instead estimates
\(k_BT\approx_{\mathrm{est}}GMm_p/R\), giving a temperature scale near
\(2\times10^7\) K with the source's solar constants. A stellar temperature
profile, composition, hydrostatic structure and stability are additional
model information; this scale balance does not establish a supernova
criterion or a quantitative nuclear reaction rate.

\section{Two-parameter regimes and a controlled approximation}

\begin{hbdefinition}{def:p15-approximation-water-model}{The linear surface-wave model}
Consider a homogeneous incompressible, inviscid, irrotational fluid of depth
\(h>0\), density \(\rho>0\), surface tension \(\gamma\geq0\) and gravity
\(g>0\). The undisturbed surface is \(z=0\) and flat bottom \(z=-h\).
For small-amplitude waves, the velocity potential satisfies
\(\phi_{xx}+\phi_{zz}=0\), \(\phi_z=0\) at the bottom, and the linear
surface conditions
\[
 \eta_t=\phi_z,\qquad
 \phi_t+g\eta-(\gamma/\rho)\eta_{xx}=0\quad(z=0).
\]
These are the stipulated linear model equations. They omit viscosity,
background flow, finite-amplitude breaking and nonflat boundaries.
\end{hbdefinition}

\begin{hbtheorem}{thm:p15-approximation-water-regimes}{proved=annex:p15.approximation.water-regimes}{The source's four corners and their error factors}
For a nonzero sinusoidal traveling mode of angular wavenumber
\(k=2\pi/\lambda>0\), this model's squared phase speed is
\[
 c^2=\left(\frac gk+\frac{\gamma k}{\rho}\right)\tanh(kh).
\]
Writing \(\ell=1/k=\lambda/(2\pi)\), the gravity/capillary and shallow/deep corners are
\[
\begin{array}{c|cc}
 &\text{gravity dominant}&\text{capillarity dominant}\\\hline
 kh\ll1&gh&\gamma h/(\rho\ell^2)\\
 kh\gg1&g\ell&\gamma/(\rho\ell)
\end{array}
\]
for \(c^2\). These are limiting forms, not exact equalities in a finite
corner. For \(t>0\),
\[
 0\leq1-\frac{\tanh t}{t}\leq\frac{t^2}{3},\qquad
 0\leq1-\tanh t\leq2e^{-2t}.
\]
If the omitted restoring term is at most \(\delta\) times the retained
one and the depth-factor relative error is at most \(\epsilon\), the
ratio of exact squared speed to its corner expression differs from one
by at most \(\delta+\epsilon+\delta\epsilon\).
\end{hbtheorem}

The source's nine-region diagram combines the four corners with formulas
retaining one or both effects. The complete expression above determines
those intervening regions within the stated linear model. Its presence is
essential: simply drawing a smooth connection between correct corners does
not prove the intervening values.

A conceptual map of classical, relativistic and quantum models likewise
needs dimensionless regime parameters. For selected speed \(v\) and action
scale \(S\), useful ratios are \(v/c\) and \(\hbar/S\). Relativistic
quantum theory, including quantum electrodynamics for electromagnetic
interactions, occupies the region where neither classical simplification
has been justified. These ratios organize modeling questions; they do not
prove a uniform classical limit for every observable, state or interaction.
Limits must state the quantities held fixed.

\begin{hbtheorem}{thm:p15-approximation-rational}{proved=annex:p15.approximation.rational}{A finite remainder and a threshold reversal}
For \(0\leq x<1\),
\[
 \frac1{1-x}=1+x+\frac{x^2}{1-x}.
\]
For \(0\leq x\leq1/2\),
\[
 1+x\leq\frac1{1-x}\leq1+x+2x^2.
\]
Thus a threshold outside that interval is decided by the corresponding
strict inequality. At \(x=0.9\), the first-order approximation is \(1.9\)
while the exact value is \(10\): comparison with \(5\) reverses.
Even at \(x=0.1\), comparing \(1+x\) with \(1.105\) gives the wrong
strict ordering of the exact function; the omitted remainder exceeds the
comparison margin.
\end{hbtheorem}

\begin{hbtheorem}{thm:p15-approximation-nonuniform}{proved=annex:p15.approximation.nonuniform}{Pointwise accuracy can fail in a moving regime}
For \(0<\varepsilon\leq1\), \(0<t\leq1\), let
\(f_\varepsilon(t)=t/(t+\varepsilon)\).
For every fixed \(t>0\), \(f_\varepsilon(t)\to1\) as
\(\varepsilon\to0^+\). Nevertheless,
\[
 \sup_{0<t\leq1}|f_\varepsilon(t)-1|=1,
 \qquad f_\varepsilon(\varepsilon)=\tfrac12.
\]
If instead \(t\geq t_0>0\), the uniform error is at most
\(\varepsilon/(t_0+\varepsilon)\). The domain restriction supplies the
uniform guarantee absent from the pointwise assertion.
\end{hbtheorem}

\begin{hbmethod}{def:p15-approximation-procedure}{A conclusion-directed approximation procedure}
Specify the target, units and decision threshold. Identify the exact model
and the dimensionless parameters controlling its regimes. Select the
rounding, characteristic-value, shape or limiting approximation and state
which relation in \cref{def:p15-approximation-relations} it establishes.
Track retained and omitted terms, including nonlinear effects of lumping.
Check the chosen regime using the resulting estimate. Obtain an error
interval when the decision needs one, and compare its margin with the
threshold. If the interval overlaps the threshold, improve the approximation
or retain the unresolved comparison. Correct easy cases can reject a bad
formula; proof, measurement or a justified error bound must supply what
those checks leave undetermined.
\end{hbmethod}

When the approximant or the model is evaluated by a program, the returned
number carries arithmetic and, for a simulation, discretization error in
addition to the approximation and input terms above.
\Cref{sec:p15-w02b-p15-074} keeps these terms separate, certifies regime
conditions on an input enclosure, and connects this procedure to
\cref{sec:p12-numerical-error-and-inference,sec:p12-simulation-results-as-evidence}.
